\subsection{Semisimplicity}

DEFINITION 7. --- Let \(\varrho\) be a unitary representation of a topological group G in a Hilbert space E. We say that \(\varrho\) is semi-simple if there exists a family \((\mathrm{F}_i)_{i\in \mathrm{I}}\) of irreducible subrepresentations of \(\varrho\) such that E is the Hilbert sum of the subspaces \(\mathrm{F}_i\) .

One sometimes also says that a semisimple unitary representation admits a discrete decomposition or is discretely decomposable.

If \((\varrho_{i})_{i\in\mathrm{I}}\) is a family of semisimple unitary representations of a topological group G, then the Hilbert sum of the representations \(\varrho_{i}\) is semisimple.

PROPOSITION 7. --- Let G be a topological group and let \(\varrho\) be a unitary representation of G in a complex Hilbert space E. The representation \(\varrho\) is semi-simple if and only if every nonzero subrepresentation of \(\varrho\) contains an irreducible subrepresentation.

Suppose that \(\varrho\) is semisimple. Let \((\mathrm{F}_i)_{i\in \mathrm{I}}\) be a family of irreducible subrepresentations of \(\varrho\) such that E is the Hilbert sum of the subspaces \(\mathrm{F}_i\) . Let F be a nonzero subrepresentation of E. There exists \(i\in \mathrm{I}\) such that the restriction to F of the orthogonal projector with image \(\mathrm{F}_i\) is not zero. This orthogonal projector then defines, by passing to subspaces, a nonzero G-morphism \(p_i\) from F into \(\mathrm{F}_i\) . By Schur's lemma, the G-morphism \(p_i\) is surjective (cor. 5 of V, p. 388). The orthogonal complement in F of the kernel of \(p_i\) is a subrepresentation of F isomorphic to \(\mathrm{F}_i\) (prop. 3 of V, p. 383), hence irreducible.

Let us prove the converse assertion. Let \(\mathcal{F}\) be the set of closed subspaces F of E that are stable under G such that the subrepresentation of \(\varrho\) in F is irreducible. Let \(\mathcal{O}\) be the set of subsets \(\mathcal{G}\) of \(\mathcal{F}\) such that the subspaces belonging to \(\mathcal{G}\) are pairwise orthogonal.

The set \(\mathcal{O}\) is of finite character (E, III, p. 34, def. 2). Let then \(\mathcal{G}\) be a maximal element of \(\mathcal{O}\) (E, III, p. 35, th. 1). Let E \(_1\) be the subspace of E that is the Hilbert sum of the irreducible subrepresentations F of \(\mathcal{G}\) . If one had E \(_1\) ≠ E, the orthogonal complement of E \(_1\) would not be zero, and the subrepresentation of G in E \(_1^{\circ}\) would by hypothesis contain an irreducible subrepresentation. Its space F would be orthogonal to the elements of \(\mathcal{G}\) , so that \(\mathcal{G} \cup \{F\} \in \mathcal{O}\) ; this contradicts the maximality of \(\mathcal{G}\) , so E = E \(_1\) , which concludes the proof.

COROLLARY 1. --- Let G be a topological group and let \(\varrho\) be a semisimple unitary representation of G in a complex Hilbert space E. Every subrepresentation of \(\varrho\) (resp. every quotient representation of \(\varrho\) ) is semisimple.

If \(\varrho_{1}\) is a subrepresentation of \(\varrho\) , then every nonzero subrepresentation of \(\varrho_{1}\) contains an irreducible subrepresentation (prop. 7), hence \(\varrho_{1}\) is semi-simple (loc. cit.).

By prop. 3 of V, p. 383, every quotient representation of \(\varrho\) is isomorphic to a subrepresentation of \(\varrho\) ; therefore it is semisimple.

COROLLARY 2. --- Every finite-dimensional unitary representation \(\varrho\) of a topological group G is semisimple.

This follows from prop. 7 and lemma 3 of V, p. 379.

COROLLARY 3. --- Let \(\varrho_{1}\) and \(\varrho_{2}\) be finite-dimensional unitary representations of a topological group G. The representations \(\varrho_{1}\) and \(\varrho_{2}\) are isomorphic if and only if \(\chi_{\varrho_{1}} = \chi_{\varrho_{2}}\) .

This follows from Corollary 2 and from A, VIII, p. 389, prop. 1, b).

